\documentclass[noauthor]{ximera}

\title{Week 11}
\author{Math 1193 Design Team}

\begin{document}
\begin{abstract}
    Solutions for Week 11.
\end{abstract}
\maketitle

\begin{problem}
Let's play a similar game to the one we played two weeks ago! 
You'll pair up to think about this problem together, but play against me this time.
We'll start with a function, 
\(f\left(x\right)=2^x\). I'll select an \(x\)-value in secret and 
tell you the corresponding output. Your job is to guess what my 
\(x\)-value was. 

\begin{enumerate}
    \item If the output is 4, what is my \(x\)-value?
    \item What if my output is 2?
    \item 1?
    \item \(\frac{1}{2}\)?
    \item 64?
    \item \(\sqrt{2}\)?
\end{enumerate}

What process did you use to find these \(x\)-values? Be as specific as possible.

If you get bored, try to play my (the author's) role.

\begin{explanation}
\begin{enumerate}
\item \(x = 2\)
\item \(x = 1\)
\item \(x = 0\)
\item \(x = -1\)
\item \(x = 6\)
\item \(x = \frac{1}{2}\)
\end{enumerate}
\end{explanation}
\end{problem}

\begin{problem}
Below is a graph of the function \(f(x) = 7^x\) with some points labeled.

\begin{center}
    \includegraphics[width=10cm]{base7.png}
\end{center}

\begin{enumerate}
    \item What is \(f\left(\frac{1}{2}\right)\)?
    \item Solve the equation \(f(x) = 49\).
    \item Does the equation \(f(x) = \frac{1}{7}\) have a solution? Can you find it?
    \item Does the equation \(f(x) = 1\) have a solution? Can you find it?
    \item Does the equation \(f(x) = -1\) have a solution? Can you find it?
    \item Does the equation \(f(x) = 2\) have a solution? Can you find it?
\end{enumerate}

\begin{explanation}
\begin{enumerate}
\item \(f\left(\frac{1}{2}\right) = \sqrt{7}\)
\item \(x = 2\)
\item Yes. Based on the graph, \(x = -1\) is a solution to \(f(x) = \frac{1}{7}\).
\item Yes. One way to tell is by seeing that there is a point on the graph with a \(y\)-coordinate
of 1. \(x = 0\) is a solution.
\item Yes. Based on the graph, the curve \(y = f(x)\) has a point with a \(y\)-coordinate
of 2. We don't have a good way to calculate it, however. It does have a name: \(\log_7(2)\).
\end{enumerate}
\end{explanation}
\end{problem}

\begin{problem}
Find a partner.

\begin{enumerate}
    \item Solve the following equations involving exponents together.
    \begin{enumerate}
        \item \(3^x = 27\)
        \item \(x^2 = 9\)
        \item \(4^x = 16\)
        \item \(4^x = 2\)
        \item \(x^4 = 2\)
        \item \(2^x = -16\)
        \item \(x^3 = -27\)
        \item \(5^x = \frac{1}{5}\)
        \item \(2^x = \frac{1}{16}\)
        \item \(9^x = \frac{1}{3}\)
        \item \(27^x = 3\)
    \end{enumerate}
    \item What do the equations have in common? What are the differences?
    \item Individually, write an equation involving an exponent, then have your partner solve it.
\end{enumerate}

\begin{explanation}
\begin{enumerate}
\item \begin{enumerate}
\item \(x = 3\)
\item \(x = -3\), \(x = 3\)
\item \(x = 2\)
\item \(x = \frac{1}{2}\)
\item \(x = -\sqrt{2}\), \(x = \sqrt{2}\)
\item no solution
\item \(x = -3\)
\item \(x = -1\)
\item \(x = -4\)
\item \(x = -\frac{1}{2}\)
\item \(x = \frac{1}{3}\)
\end{enumerate}
\item Some equations involve exponential expressions, where the 
variable is an exponent, and some involve polynomials, where the variable
is raised to some fixed power. The equations have 0, 1, or more solutions.
\item An example would be \(3^x = -3\), which does not have a solution.
\end{enumerate}
\end{explanation}
\end{problem}

\begin{problem}
We'll investigate powers of 5. Below is a graph of the function \(f(x) = 5^x\).

\begin{center}
    \includegraphics[width=10cm]{expGrowth.png}
\end{center}

\begin{enumerate}
    \item Solve \(5^x = 25\). What point on the graph of \(f\) does this solution correspond to?
    \item Solve \(5^x = 125\). What point on the graph of \(f\) does this solution correspond to?
    \item Solve \(5^x = 625\). What point on the graph of \(f\) does this solution correspond to?
    \item Consider the equation \(5^x = -5\). Does this correspond to a point on the graph of \(f\)?
    What does this mean about whether there are solutions to the equation?
    \item Consider the equation \(5^x = 50\). Does this correspond to a point on the graph of \(f\)?
    What does this mean about whether there are solutions to the equation?
    \item Now we'll introduce a mystery number \(a\). The
    definition of \(a\) is that \(5^a = 50\). That is, \(a\)
    is the solution of the equation \(5^x = 50\).
    \begin{enumerate}
        \item Is \(a\) a whole number? (Whole numbers are 0, 1, 2, 3, 4, ...)
        \item What two consecutive integers does \(a\) fall between? (Some 
        examples of consecutive integers are \(-2\), \(-1\); or \(5\), \(6\); or \(101\), \(102\).)
    \end{enumerate}
\end{enumerate}

\begin{explanation}
\begin{enumerate}
    \item \(x = 2\). This corresponds to the point \((2, 25)\).
    \item \(x = 3\). This corresponds to the point \((3, 125)\).
    \item \(x = 4\). This corresponds to the point \((4, 625)\).
    \item There is no point with a 
    \(y\)-coordinate of -5 on the curve, so there are no solutions
    to \(5^x = -5\).
    \item There is a point with \(y\)-coordinate
    50 on the curve, so there is a solution to \(5^x = 50\).
    \item \begin{enumerate}
        \item \(a\) is not a whole number.
        \item \(a\) lies between 2 and 3, since 
        \(25 < 50 < 125\) and therefore, by rewriting, \(5^2 < 5^a < 5^3\).
        Since exponentials are increasing functions, \(a\) must be some number between 2 and 3.
    \end{enumerate}
\end{enumerate}
\end{explanation}
\end{problem}

\end{document}
